\section{Executive Accountability：责任共享，但不能稀释}

重大安全事件往往暴露组织结构：谁收到事实、谁能要求更多调查、谁批准付款、谁决定不通知、谁向董事会报告。把责任全部压给 CSO 会诱发防御性文化；把责任平均分散给十几个委员会，又会让任何人都不真正 accountable。成熟治理需要 shared responsibility 与 named decision owner 同时存在。

\subsection{CSO、GC、CEO 与 board 的 decision rights}

CSO 应确保风险被准确升级，GC 应确保法律问题由有资质团队处理，CEO 应为资源和重大经营选择负责，board 或 disclosure committee 应提供监督与 challenge。CSO 不能以“legal approved”替代自己的技术真实性义务，GC 也不能在缺少 evidence 时单独判断影响，CEO 更不能只在公开披露前才第一次听说安全架构。

\lecturefigure{12-executive-accountability.png}{网络安全责任由 CSO、GC、CEO 与 board 共同承担，但每项决定仍需明确 owner。}{本地字幕 25:20--35:20；治理概念重绘。}

\begin{knowledgebox}{读图：共享信息，不共享模糊责任}
CSO 提供 technical scope、security recommendation 与 escalation facts；GC 提供 legal duties、regulator strategy 与 privilege advice；CEO 作 enterprise risk 和 disclosure decision；board/committee 监督、挑战并要求复盘。图中边界不是 silo：每个角色都能 challenge 其他 workstream，但不能把自己的 deliverable 推给别人。最终记录应能回答谁负责事实、谁负责分析、谁批准行动。
\end{knowledgebox}

\teachervoice{讲者建议安全领导者在危机前就与 CEO、GC 建立关系，而不是第一次重大通话就要求他们理解复杂 incident。课堂提醒：executive readiness 是一种预先训练的 interface，需要定期 brief、tabletop 和明确的“何时叫醒我”阈值。}

\begin{warningbox}{“我已经告诉过他们”不是治理完成}
单次邮件或口头提醒不能证明接收者理解风险、拥有 decision right 或完成后续行动。高风险升级应包含 severity、evidence、unknowns、recommended actions、deadline、明确 ask 与确认；若决策被延迟或拒绝，应记录 owner、理由和下一次复审。
\end{warningbox}

\begin{importantbox}{保护安全领导者的正确方式}
保护不等于制造自利 memo。组织应提供清晰 job authority、written escalation、independent legal advice、董事会 access、D\&O/indemnification 与适当 employment support；领导者则应保持事实诚实、避免越权作法律结论、拒绝删除不利记录，并在重大分歧时使用正式升级渠道。
\end{importantbox}

\subsection{本章小结}

Accountability 的目标是让有能力、有信息且有权的人作决定。CSO 不应成为所有安全风险的单点承担者，也不能因“集体决策”而放弃准确升级和专业判断。

